% GENERATED by md2tex.py from ../draft_v5_en/body/10_operational_validation.md -- do not edit by hand.
% Change the Markdown and run `python md2tex.py`; edits here are overwritten.

\chapter{Operational validation}
\label{ch:10}

\section{Purpose of the chapter}
\label{sec:10.1}

The earlier chapters produce measurements. This chapter turns them into a validation statement that includes its own limits. It is the first half of adaptation A5 in practice: the conclusion is no longer \enquote{the system is safe} but a bounded claim about which requirements are met, in which domain, with what evidence, and with which residual risks stated.

\section{The evidence}
\label{sec:10.2}

The statement is based on three sets of evidence. All of them have version-controlled artefacts, and every run records its reproducibility metadata: code revision, cage configuration identifier, checkpoint identifier, seed and timestamp.

The perfect perception arm gives 1,260 runs with a global verdict of satisfied. It works as a control, because it shows what happens when perception is not a problem. The camera arm gives the scenario campaigns on a policy that learned from pixels, and ends with the reference campaign of 1,890 runs without errors that supports the verdict of this chapter. The study of variation between seeds shows that behaviour is not the same when the training procedure is repeated.

\textbf{Physical evidence exists, but it comes from bring-up and is not verdict evidence.} The chain has run on the real vehicle and has driven 18.05 m of the circuit in one segment without triggering any safety rule (Chapter~\ref{ch:09}). However, no run was done under the scenario library protocol, in enforcement mode, and with the perception contract used to score the campaigns. For that reason the physical column of the table below is marked not executed.\sidenote{Not executed rather than pending: no measurement is in progress, and no data taken outside the protocol is used to fill the column.} This is more precise than saying that there is no data. Having driven is not the same as having scored, and filling the column with measurements taken outside the protocol would be exactly the kind of claim the framework is meant to prevent. The limits this imposes are stated in \S\ref{sec:10.4}d–\allowbreak{}e.

\section{Final verdict table}
\label{sec:10.3}

% layout: table*, fixed, \small, 164 of 165 mm
\begin{table*}[htbp]
\caption[Final verdicts per requirement.]{Final verdicts per requirement. The physical column is not scored in this work: no run on hardware was done under the scenario protocol, and all of them were done in monitoring mode (\S\ref{sec:10.4}e). It is marked \enquote{not executed} and not \enquote{pending}, so as not to suggest a measurement in progress.}\label{tab:10.1}
\small\setlength{\tabcolsep}{4pt}
\begin{tabular}{@{}>{\RaggedRight\arraybackslash}p{41.2mm}>{\centering\arraybackslash}p{10.6mm}>{\RaggedRight\arraybackslash}p{25.9mm}>{\RaggedRight\arraybackslash}p{49.7mm}>{\RaggedRight\arraybackslash}p{24.9mm}@{}}
  \toprule
  \tabhead{Requirement} & \tabhead{Class} & \tabhead{Control arm} & \tabhead{\textbf{Camera arm (reference campaign)}} & \tabhead{Physical \emph{(out of scope — \S\ref{sec:12.4} T2)}} \\
  \midrule
  \sr{001} lateral deviation & A & Satisfied & Satisfied & not executed \\
  \addlinespace[2pt]
  \sr{002} heading stability & A & Satisfied & Literal: fail (recovery clause); own criterion: satisfied & not executed \\
  \addlinespace[2pt]
  \sr{003} predictive time to departure & A & Satisfied & Literal: fail (same clause); own criterion: satisfied & not executed \\
  \addlinespace[2pt]
  \sr{004} speed ceiling & A & Satisfied & Satisfied (rule not activated in the campaign; see \S\ref{sec:10.4}b and \S\ref{sec:10.4}f) & not executed \\
  \addlinespace[2pt]
  \sr{005} emergency stop & A & Satisfied & Satisfied & not executed \\
  \addlinespace[2pt]
  \sr{006} actuation smoothness & B & Satisfied (own metric) & Satisfied (840/840 in enforcement) & not executed \\
  \addlinespace[2pt]
  \sr{007} state validity & A & Satisfied & Satisfied & not executed \\
  \addlinespace[2pt]
  \sr{008} external stop & A & Satisfied & Satisfied & not executed \\
  \addlinespace[2pt]
  \sr{009} \emph{liveness} & B & Documented abstention & Satisfied (own measurements) & not executed \\
  \addlinespace[2pt]
  \sr{010} rule composition & B & Documented abstention & Not satisfied — finding, does not block & not executed \\
  \addlinespace[2pt]
  \sr{011} heading variance & B & Satisfied & Literal: fail (same inherited clause); own criterion: satisfied (3.77° < 5°) & not executed \\
  \addlinespace[2pt]
  \sr{012} following under a degraded camera & A & n/a & Satisfied & not executed \\
  \addlinespace[2pt]
  \sr{013} safe perception degradation & A & n/a & Satisfied & not executed \\
  \addlinespace[2pt]
  \sr{014} estimator plausibility & A & n/a & Satisfied & not executed \\
  \bottomrule
\end{tabular}
\end{table*}

\begin{figure*}[htbp]
  \includegraphics[width=0.84\widefigurewidth]{fig_10_1_traceability_case_sr001.png}
  \caption[The traceability chain instantiated end to end.]{One row of Table~\ref{tab:10.1} followed backwards to the evidence behind it: the chain \texttt{Hazard → Requirement → Rule → Scenario → Metric → Evidence → Verdict}, shown for the most important requirement. Every link is an identifier in a versioned artefact, and \texttt{check\_\allowbreak{}traceability.\allowbreak{}py} blocks the review gate if an orphan appears in either direction. This is how a cell of the table becomes auditable. The evidence is the reference campaign (\texttt{campaign\_\allowbreak{}2d\_\allowbreak{}ppo550k}, \texttt{2-D PPO 550k}; Table~\ref{tab:7.1}).}\label{fig:10.1}
\end{figure*}

Seen in groups: thirteen of the fourteen requirements are satisfied on their documented criterion. One is not, and it is class B, so it does not block the verdict; it is reported as it is, without being explained away. None is left without a verdict, left out of the table, or pending in the simulation column. The coverage criterion from the framework evaluation, which asked for a verdict on 100\,\% of the requirements even when the verdict is uncomfortable, is met, including that uncomfortable verdict.

Two groups of cells need a comment. The three requirements with a literal failure (\sr{002} and \sr{003} in class A, \sr{011} in class B) all fail through the same inherited clause of \S\ref{sec:8.3.2}; they keep that failure in the record next to the explanation, and are not rewritten. And the speed ceiling is marked satisfied with an important condition, explained below: it is satisfied without ever having been tested.

\section{Bounded validation statement}
\label{sec:10.4}

\begin{quote}
\textbf{Statement.} Within the specified operational domain, on a single circuit in simulation, with the two-dimensional camera policy chosen by closed-loop evaluation and with the cage configuration identified by its \emph{hash}, the system meets thirteen of its fourteen safety requirements on their documented criterion. The literal global verdict of the campaign is \verdictnot{}, and it is caused entirely by one performance clause, the heading recovery time, in one scenario, without any class A safety condition being broken.

\textbf{Inside the operational domain and with the cage active, there is no contact with the road edge} over the 555 enforcement runs that start inside it. The same policy without the cage makes sixty over those same 555. The claim is bounded to the domain on purpose: the 390 enforcement runs that start outside it still contain 56 contacts (Table~\ref{tab:8.3}). The requirement that is not met is the one on consistent rule composition when rules fire at the same time. It is class B, it was measured on two different policies, it was reduced by half with better training, and it stays the same in kind. It is stated as a design limitation of the envelope and as future work.
\end{quote}

\textbf{Explicit limits of this statement.} They are listed here, as part of the statement and not in a separate section, because each one limits a specific claim made above.

\textbf{(a) Scope.} The statement is valid in simulation, on one circuit, and with one seed in the verdict campaign. Variation between seeds is studied separately and is not uniform.

\textbf{(b) What was not tested.} The cage's speed ceiling did not fire once in the 1,890 runs, because the policy's operating speed stays below the lower limit of that envelope. The requirement is satisfied trivially\sidenote{Trivially satisfied: the requirement was never challenged, because the speed never reached the rule's threshold, so its rule was never exercised.}, and its rule has never been tested from above. This is a limitation of the operating point, not of the specification, and it is stated so that it is not read as evidence of robustness.

\textbf{(c) The measured dependency.} For this policy, staying in the lane on the tightest curves depends on the rate limiter. The dependency is measured, and its origin is inferred.

\textbf{(d) Transfer.} Physical evidence exists but is not scored, so no claim in this statement applies to the real platform. Chapter~\ref{ch:09}, and not this section, makes the one claim that can be made: the policy behind this verdict does not transfer to the real vehicle, and the one that does transfer is a later retraining. The system validated in this statement and the system that drives on hardware are not the same.

\textbf{(e) What the hardware has not tested at all, which is the biggest limit of this work.} All physical runs were done in monitoring mode. In that mode the cage checks, publishes and logs its rules, but it does not change the action, so the safe action is identical to the policy's raw action in every logged cycle. The envelope has therefore never acted on the real vehicle. The physical step measures what the cage \emph{would have said}, plus the availability cost of its emergency latch. It does not measure how the cage changes the trajectory of a physical vehicle, which is exactly the property this statement certifies in simulation. This is a limit of scope and not a negative result. There is no physical evidence against the envelope and none in its favour, and there will be none until a scored run is done in enforcement mode.

\textbf{(f) One rule that the hardware did test.} The speed ceiling is marked satisfied here without having fired once, and \S\ref{sec:9.3.5} adds two physical measurements that limit this cell without re-scoring it: in the deployed configuration the rule cannot fire on commanded motion, and the only times it did fire were caused by velocity errors from the pose sensor. For this statement that means one specific and limited thing: the verdict on this requirement depends entirely on the fact that its trigger condition was never reached in simulation, and not on evidence that the rule acts correctly when it is reached. It is the requirement with the weakest support in the table, and it is marked as such.

\section{Relation to the methodological thesis}
\label{sec:10.5}

The statement above is what the framework promised to produce: not a yes-or-no judgement, but a bounded statement that can be traced back to its evidence and that states clearly what it does not cover. Its form matters as much as its content, because that is what separates an honest validation from an approval. Every claim can be followed back to a set of logged runs, and every limit is written inside the statement rather than in a separate section a reader in a hurry could skip. The three unfavourable results — a negative global verdict, an unmet requirement and a rule that was never tested — are all inside it, and Chapter~\ref{ch:11} evaluates what that says about the framework.
